\section{Колоночная сетка и изменение геометрии}
\label{sec:mesh}

\draft{Перед построением сетки полилинии контактов могут упрощаться
алгоритмом Рамера--Дугласа--Пекера (RDP)~\cite{ramer1972,douglas1973}
с заданным геометрическим допуском. Уменьшение числа вершин позволяет
сократить число колонок. Влияние этой операции на точность T и Ro
и время расчёта исследуется в разделе~\ref{sec:performance}.}

\draft{Область разбивается на колонки прямоугольных ячеек с независимыми
ограничениями шагов $\Delta x$ и $\Delta z$. Уровни соседних колонок могут
не совпадать; их общая вертикальная граница разбивается на подграни.
Рисунок~\ref{fig:mesh} показывает общую сетку и участок у разлома.}
\reportfigure{02_column_mesh.png}
  {Современная сетка обратного набора при $\Delta x=200$ м, $\Delta z=50$ м
  и её увеличение у разлома. Цвета обозначают литотипы.}
  {fig:mesh}

\draft{Ячейки нового слоя создаются сразу на всю эпоху и раскрываются снизу
вверх по мере отложения. При конечной мощности $h_{\mathrm{end}}$ и относительной
глубине $a\in[0,1]$ принята следующая интерполяция:}
\begin{gather}
 \theta=\frac{t-t_k}{t_{k+1}-t_k},\qquad h(x,t)=\theta h_{\mathrm{end}}(x),\\
 z(a,t)=z_{\mathrm{surf}}(x,t)
          +h_{\mathrm{end}}(x)\max(0,\theta+a-1),\qquad
 t_{\mathrm{dep}}=t_k+(1-a)(t_{k+1}-t_k).
 \label{eq:deposition}
\end{gather}
\draft{Численно нулевая высота заменяется $\varepsilon=0{,}05$ м;
суммарная мощность $N$ сжатых ячеек равна $N\varepsilon$.
Внешний контур сохраняется перераспределением толщин. На рис.~\ref{fig:growth}
показаны раскрытие ячеек, моменты отложения и увеличение сжатого участка.}
\reportfigure{03_deposition_growth.png}
  {Новый слой в начале, середине и конце эпохи.
  T хранится в центрах ячеек, Ro --- на горизонтальных границах;
  справа указаны моменты отложения относительно начала эпохи.}
  {fig:growth}

\draft{В каждой колонке с $N$ ячейками зрелость и интегралы хранятся в $N+1$
точке. Для расчёта Ro температура восстанавливается линейно по вертикали.
При смене сетки T и интегралы интерполируются внутри соответствующих слоёв;
перенос не объявляется консервативным.}
